\chapter{Programas elaborados}
\label{ap:codigo}

La simulación completa (configuración de modelos, evaluación de los cuatro métodos en Fase~1 y Fase~2, y generación de capturas y figuras) se implementa en MATLAB/Simulink. La red profunda se representa como una política de inferencia matricial entrenada \emph{offline}, el Q-Learning como una tabla de valores estado--acción y el supervisor como funciones MATLAB llamadas desde los bloques de control. El mismo banco fija los parámetros de planta, perturbaciones, límites de actuación y escenarios usados en los capítulos de resultados.

Se incluyen extractos del código que implementan la dinámica de la planta, los controladores, la actualización tabular y el supervisor.

En los modelos Simulink generados, las politicas no quedan embebidas como bloques decorativos: se invocan funciones MATLAB externas. Para el servo se usan \texttt{act08\_servo\_fuzzy}, \texttt{act08\_servo\_dl} y \texttt{act08\_servo\_ql}; para la celda R2ET, \texttt{act08\_r2et\_fuzzy}, \texttt{act08\_r2et\_dl} y \texttt{act08\_r2et\_ql}; y para la capa de seguridad, \texttt{act08\_supervisor\_fuzzy}. El selector de metodo del modelo permite elegir la politica sin cambiar la planta ni el supervisor.

\lstset{
  language=Matlab,
  basicstyle=\ttfamily\scriptsize,
  numbers=left, numberstyle=\tiny,
  breaklines=true, columns=fullflexible,
  keepspaces=true, upquote=true,
  frame=single, tabsize=4,
  keywordstyle=\color{UniPrimaryDeep}\bfseries,
  commentstyle=\color{UniMuted}\itshape,
  stringstyle=\color{UniPrimary!80!black},
  showstringspaces=false
}

\begin{lstlisting}[language=Matlab, caption={Modelo servo reducido del Grupo~1 y controlador PID con anti-windup.}]
% Modelo reducido de primer orden (ZOH, Ts = tau/12) - ecuacion 2.3
A  = 0.92;     % polo discreto
B  = 0.08;     % ganancia de control
Bd = 0.024;    % ganancia de perturbacion de carga
r  = 0.60;     % referencia de operacion normalizada
u_ff = r*(1 - A)/B;

% y(k+1) = A*y(k) + B*u(k) - Bd*d(k)

function [u,estado] = pid_servo(e,estado,Kp,Ki,Kd,u_ff)
    de = e - estado.e_ant;
    u_raw = u_ff + Kp*e + Ki*estado.I + Kd*de;
    u = min(max(u_raw,0.05),1.0);

    satura_alta = (u_raw > 1.0) && (e > 0);
    satura_baja = (u_raw < 0.05) && (e < 0);
    if ~(satura_alta || satura_baja)
        estado.I = min(max(estado.I + e,-6),6);
    end
    estado.e_ant = e;
end
\end{lstlisting}

\begin{lstlisting}[language=Matlab, caption={Controlador difuso: inferencia AND-mínimo y defuzzificación por singletons.}]
function u = velocidad_difusa(mu_e,mu_de,mu_d,u_base)
% FAM de 35 reglas: 7 conjuntos de error x 5 conjuntos de derivada.
% Consecuente: c = u_base + nivel_e*0.12 + nivel_de*0.05

    nivel_e  = [-3 -2 -1 0 1 2 3];
    nivel_de = [-2 -1 0 1 2];
    num = 0;
    den = 0;

    for ie = 1:numel(nivel_e)
        for ide = 1:numel(nivel_de)
            w = min(mu_e(ie),mu_de(ide));   % AND = minimo
            c = min(max(u_base + 0.12*nivel_e(ie) + ...
                        0.05*nivel_de(ide),0.05),0.98);
            num = num + w*c;
            den = den + w;
        end
    end

    % Regla de perturbacion: SI d = CON ENTONCES velocidad alta
    w_d = mu_d(2);
    c_d = min(u_base + 0.30,0.98);
    num = num + w_d*c_d;
    den = den + w_d;

    if den > 1e-9
        u = min(max(num/den,0.05),0.98);
    else
        u = u_base;
    end
end
\end{lstlisting}

\begin{lstlisting}[language=Matlab, caption={Actualización TD(0) del Q-Learning y codificación de estado para la coordinación de esteras.}]
% Estado: (nivel_E1_bin, nivel_E2_bin, perturbacion_bin) -> 9x9x2 = 162 estados.
edges = [0.0 0.5 0.9 1.2 1.4 1.6 1.8 2.1 2.5 3.1];

acciones = [
    0.30 0.50 1.00
    0.46 0.50 1.00
    0.58 0.50 1.00
    0.70 0.50 0.95
    0.84 0.50 0.85
    0.95 0.50 0.78
    0.52 0.30 1.00
    0.52 0.70 1.00
    0.72 0.30 0.90
    0.72 0.70 0.90
    0.62 0.50 1.00
];

% Regla TD(0): actualizacion de Bellman
r = -1.5*iae - 30.0*violacion - 8.0*alarma - 0.20*desbalance;
Q(s,a) = Q(s,a) + alpha*(r + gamma*max(Q(s_siguiente,:)) - Q(s,a));
\end{lstlisting}

\begin{lstlisting}[language=Matlab, caption={Bucle de simulación Fase~2: controlador difuso más supervisor de capacidad.}]
for k = 2:N_muestras
    lambda = demanda(k);
    d = perturbacion(k);
    perdida = [0.35*d, 0.175*d];
    nmax = max(n(k-1,:));
    e = n(k-1,:) - n_objetivo;
    de = e - e_anterior;

    % 1. Controlador difuso: velocidad por estera
    u = zeros(1,2);
    for i = 1:2
        u(i) = velocidad_difusa(mu_e{i},mu_de{i},mu_d,u_base);
    end

    % 2. Supervisor de capacidad
    if nmax > umbral_alarma
        u = min(max(u + 0.14,0.05),1.0);
        admision = 0.72;
    elseif d > 0
        u = min(max(u + 0.14,0.05),1.0);
        admision = 0.86;
    elseif nmax > umbral_riesgo
        u = min(max(u + 0.07,0.05),1.0);
        admision = 0.90;
    end

    % 3. Reparto del robot y balance de masa
    r1 = reparto_difuso(n(k-1,1) - n(k-1,2));
    entrada = lambda*admision*[r1,1-r1];
    salida = eta*u.*(1 - perdida);
    n(k,:) = max(0,n(k-1,:) + entrada - salida);
end
\end{lstlisting}


\begin{lstlisting}[language=Matlab,caption={Configuración del sistema difuso para el controlador inteligente directo},label={lst:fis_control_directo}]
% Sugeno de orden cero: antecedentes Mamdani (AND-min) + singletons + promedio ponderado
fis = sugfis('Name','fis_control_directo');

% Entrada e: 7 conjuntos, universo documentado en la Tabla 4.2
fis = addInput(fis,[-1.5 1.5],'Name','e');
e_names = {'NL','NM','NS','ZE','PS','PM','PL'};
e_types = {'trapmf','trimf','trimf','trimf','trimf','trimf','trapmf'};
e_params = {[-1.5 -1.5 -0.7 -0.4],[-0.6 -0.35 -0.12],[-0.22 -0.11 0], ...
            [-0.07 0 0.07],[0 0.11 0.22],[0.12 0.35 0.6],[0.4 0.7 1.5 1.5]};
for i = 1:numel(e_names)
    fis = addMF(fis,'e',e_types{i},e_params{i},'Name',e_names{i});
end

% Entrada de: 5 conjuntos, universo documentado en la Tabla 4.2
fis = addInput(fis,[-0.5 0.5],'Name','de');
de_names = {'NB','NS','ZE','PS','PB'};
de_types = {'trapmf','trimf','trimf','trimf','trapmf'};
de_params = {[-0.5 -0.5 -0.22 -0.10],[-0.18 -0.08 0],[-0.05 0 0.05], ...
             [0 0.08 0.18],[0.10 0.22 0.5 0.5]};
for i = 1:numel(de_names)
    fis = addMF(fis,'de',de_types{i},de_params{i},'Name',de_names{i});
end

% Entrada d: SIN/CON
fis = addInput(fis,[0 1],'Name','d');
fis = addMF(fis,'d','trapmf',[0 0 0.35 0.55],'Name','SIN');
fis = addMF(fis,'d','trapmf',[0.35 0.55 1 1],'Name','CON');

% Salida u_star con consecuentes singleton
fis = addOutput(fis,[0.05 0.98],'Name','u_star');
u_base = 0.60;
e_level = [-3 -2 -1 0 1 2 3];
de_level = [-2 -1 0 1 2];
rules = [];
out_idx = 1;
for ie = 1:numel(e_level)
    for ide = 1:numel(de_level)
        c = min(max(u_base + 0.12*e_level(ie) + 0.05*de_level(ide),0.05),0.98);
        fis = addMF(fis,'u_star','constant',c,'Name',sprintf('C%02d',out_idx));
        rules = [rules; ie ide 0 out_idx 1 1];  % 0 = no importa
        out_idx = out_idx + 1;
    end
end
fis = addMF(fis,'u_star','constant',min(u_base + 0.30,0.98),'Name','PERT');
rules = [rules; 0 0 2 out_idx 1 1];            % regla d=CON -> velocidad alta

% Operadores de inferencia
fis.AndMethod = 'min';
fis.OrMethod = 'max';
fis.DefuzzificationMethod = 'wtaver';

fis = addRule(fis,rules);
fis_control_directo = fis;
\end{lstlisting}

La función siguiente implementa únicamente la barrera dura de capacidad usada en el modelo de coordinación del Capítulo~\ref{ch:arquitectura} (Figura~5.16); los refuerzos de velocidad y los regímenes de riesgo/alarma de las ecuaciones del Capítulo~\ref{ch:supervisor} se implementan en \texttt{supervisor\_inteligente} (Listing~\ref{lst:supervisor_inteligente}) dentro del modelo supervisor del Capítulo~\ref{ch:supervisor}.

\begin{lstlisting}[language=Matlab,caption={Función del supervisor de capacidad para la coordinación R2ET},label={lst:supervisor_capacidad}]
function [u1,u2,r1,alpha] = supervisor_capacidad(u1s,u2s,r1s,n1,n2,nmax)

% Senales propuestas por el controlador primario
u1 = u1s;
u2 = u2s;
r1 = r1s;
alpha = 1;

margen = 0.05;

llena1 = n1 >= (nmax - margen);
llena2 = n2 >= (nmax - margen);

% Bloqueo de alimentacion hacia esteras llenas
if llena1 && ~llena2
    r1 = 0.20;
end

if llena2 && ~llena1
    r1 = 0.80;
end

if llena1 && llena2
    alpha = 0.72;
    r1 = 0.5;
end

% Limites fisicos
if r1 < 0.20
    r1 = 0.20;
elseif r1 > 0.80
    r1 = 0.80;
end

if u1 < 0.05
    u1 = 0.05;
elseif u1 > 1
    u1 = 1;
end

if u2 < 0.05
    u2 = 0.05;
elseif u2 > 1
    u2 = 1;
end

end
\end{lstlisting}

\begin{lstlisting}[language=Matlab,caption={Función del supervisor inteligente para ajuste de ganancias, caudal y reparto},label={lst:supervisor_inteligente}]
function [Kp,Ki,Kd,u_plus,alpha,r1,modo] = supervisor_inteligente(n1,n2,d,theta,nmax_lim)

nmax = max(n1,n2);

umbral_riesgo = 0.80*nmax_lim;
umbral_alarma = 0.90*nmax_lim;
umbral_lleno  = nmax_lim - 0.05;

% Valores por defecto: modo normal
Kp = 1.0;
Ki = 1.0;
Kd = 0.0;
u_plus = 0.0;
alpha = 1.0;
modo = 0;

% Seleccion del regimen de operacion
if nmax > umbral_alarma
    Kp = 1.55;
    Ki = 0.65;
    Kd = 0.0;
    u_plus = 0.14;
    alpha = 0.72;
    modo = 2;

elseif d > 0
    Kp = 1.55;
    Ki = 0.65;
    Kd = 0.0;
    u_plus = 0.14;
    alpha = 0.86;
    modo = 2;

elseif nmax > umbral_riesgo
    Kp = 1.25;
    Ki = 0.82;
    Kd = 0.0;
    u_plus = 0.07;
    alpha = 0.90;
    modo = 1;

else
    Kp = 1.0;
    Ki = 1.0;
    Kd = 0.0;
    u_plus = 0.0;
    alpha = 1.0;
    modo = 0;
end

% Reparto base del robot
r1 = 0.5;

if abs(theta) > 0.15
    r1 = 0.5 + 0.8*theta;
end

% Barrera dura de capacidad
if n1 >= umbral_lleno && n2 < umbral_lleno
    r1 = 0.20;
end

if n2 >= umbral_lleno && n1 < umbral_lleno
    r1 = 0.80;
end

if n1 >= umbral_lleno && n2 >= umbral_lleno
    alpha = 0.72;
    r1 = 0.5;
end

if r1 < 0.20
    r1 = 0.20;
elseif r1 > 0.80
    r1 = 0.80;
end

end
\end{lstlisting}

\begin{lstlisting}[language=Matlab,caption={Función del PID local de ocupación para las esteras R2ET},label={lst:pid_ocupacion}]
function [u1,u2,e1,e2] = pid_ocupacion(n_ref,n1,n2,Kp,Ki,Kd,u_plus)

persistent I1 I2 e1_ant e2_ant

if isempty(I1)
    I1 = 0;
    I2 = 0;
    e1_ant = 0;
    e2_ant = 0;
end

Ts = 1;

% Error de ocupacion respecto a la referencia
e1 = n1 - n_ref;
e2 = n2 - n_ref;

I1 = I1 + e1*Ts;
I2 = I2 + e2*Ts;

de1 = (e1 - e1_ant)/Ts;
de2 = (e2 - e2_ant)/Ts;

u1_raw = Kp*e1 + Ki*I1 + Kd*de1 + u_plus;
u2_raw = Kp*e2 + Ki*I2 + Kd*de2 + u_plus;

u1 = u1_raw;
u2 = u2_raw;

if u1 < 0.05
    u1 = 0.05;
elseif u1 > 1
    u1 = 1;
end

if u2 < 0.05
    u2 = 0.05;
elseif u2 > 1
    u2 = 1;
end

e1_ant = e1;
e2_ant = e2;

end
\end{lstlisting}
